% from ega2-1-fr.tex
\oldpage[II]{58}

où la graduation du second membre provient naturellement de celle de
$\mathcal{S}$. On peut en effet se borner au cas où
$Y=\operatorname{Spec}(A)$ est affine, $\mathcal{M}=\widetilde{M}$ et
$\mathcal{S}=\widetilde{S}$, la $A$-algèbre graduée $S$ étant engendrée par
$S_1$, de sorte que lorsque $f$ parcourt $S_1$, les $D_+(f)$ forment un
recouvrement de $X$. Revenant aux définitions
(\hyperref[II.2.6.2-fr]{2.6.2}) on voit alors, compte tenu de
(\hyperref[I.1.6.7-fr]{I, 1.6.7}), que la restriction à $D_+(f)$ de
l'homomorphisme (\hyperref[II.3.3.2.4-fr]{3.3.2.4}) correspond
(\hyperref[I.1.3.8-fr]{I, 1.3.8}) à l'homomorphisme de $S_{(f)}$-modules
$M_n\otimes_A S_{(f)}\to(S(n))_{(f)}$ qui, à $x\otimes1$ ($x\in M_n$), fait
correspondre $x/1$~; cela démontre notre assertion.
\end{env}

\begin{proposition}[3.3.3]
\label{II.3.3.3-fr}
Pour toute section $f\in\Gamma(Y,\mathcal{S}_d)$ ($d>0$), $X_f$ est identique
à l'ensemble des points de $X$ où $\alpha_d(f)$ (considérée comme section de
$\mathcal{O}_X(d)$) ne s'annule pas
(\hyperref[0.5.5.2-fr]{0, 5.5.2}).
\end{proposition}

\begin{proof}
($\alpha_d(f)$ est une section de $p_*(\mathcal{O}_X(d))$ au-dessus de $Y$,
mais par définition une telle section est aussi une section de
$\mathcal{O}_X(d)$ au-dessus de $X$
(\hyperref[0.4.2.1-fr]{0, 4.2.1})). La définition de $X_f$
(\hyperref[II.3.1.4-fr]{3.1.4}) ramène au cas où $Y$ est affine, qui a été
traité dans (\hyperref[II.2.6.3-fr]{2.6.3}).
\end{proof}

\begin{env}[3.3.4]
\label{II.3.3.4-fr}
Nous supposerons désormais, outre l'hypothèse du début de ce numéro, que pour
tout $\mathcal{O}_X$-Module quasi-cohérent $\mathcal{F}$, les
$p_*(\mathcal{F}(n))$ sont \emph{quasi-cohérents} sur $Y$, et par suite
$\Gamma_*(\mathcal{F})=\bigoplus_{n\in\mathbf{Z}}p_*(\mathcal{F}(n))$ est
aussi un $\mathcal{O}_Y$-Module quasi-cohérent
(\hyperref[I.1.4.1-fr]{I, 1.4.1} et
\hyperref[I.1.3.9-fr]{I, 1.3.9})~; cette circonstance se produira toujours en
particulier si $X$ est \emph{de type fini} sur $Y$
(\hyperref[I.9.2.2-fr]{I, 9.2.2}). On en conclut que
$(\Gamma_*(\mathcal{F}))^{\sim}$ est défini et est un
$\mathcal{O}_X$-Module quasi-cohérent. Pour tout ouvert affine $U$ de $Y$, on
a (\hyperref[I.1.3.9-fr]{I, 1.3.9} et
\hyperref[II.2.5.4-fr]{2.5.4})
\begin{align*}
  \left(\Gamma\left(U,\bigoplus_{n\in\mathbf{Z}}
  p_*(\mathcal{F}(n))\right)\right)^{\sim}
  &=\bigoplus_{n\in\mathbf{Z}}
  \bigl(\Gamma(U,p_*(\mathcal{F}(n)))\bigr)^{\sim}\\
  &=\bigoplus_{n\in\mathbf{Z}}
  \bigl(\Gamma(p^{-1}(U),\mathcal{F}(n))\bigr)^{\sim}\\
  &=\left(\bigoplus_{n\in\mathbf{Z}}
  \Gamma(p^{-1}(U),\mathcal{F}(n))\right)^{\sim}\\
  &=\bigl(\Gamma_*(\mathcal{F}|p^{-1}(U))\bigr)^{\sim}
\end{align*}
et par suite (\hyperref[II.2.6.4-fr]{2.6.4}) on a un homomorphisme canonique
\[
  \beta_U:\left(\Gamma\left(U,\bigoplus_{n\in\mathbf{Z}}
  p_*(\mathcal{F}(n))\right)\right)^{\sim}
  \longrightarrow\mathcal{F}|p^{-1}(U).
\]

En outre, la commutativité de
(\hyperref[II.2.8.13.2-fr]{2.8.13.2}) montre que ces homomorphismes commutent
avec les opérations de restriction sur $Y$~; on en déduit donc un
homomorphisme canonique fonctoriel
\[
\label{II.3.3.4.1-fr}
  \beta:(\Gamma_*(\mathcal{F}))^{\sim}\longrightarrow\mathcal{F}
\tag{3.3.4.1}
\]
(aussi noté $\beta_{\mathcal{F}}$) pour les $\mathcal{O}_X$-Modules
quasi-cohérents.
\end{env}

\begin{proposition}[3.3.5]
\label{II.3.3.5-fr}
Soient $\mathcal{M}$ un $\mathcal{S}$-Module gradué quasi-cohérent,
$\mathcal{F}$ un $\mathcal{O}_X$-Module quasi-cohérent~; les homomorphismes
composés
\[
\label{II.3.3.5.1-fr}
  \widetilde{\mathcal{M}}
  \xrightarrow{\widetilde{\alpha}}
  (\Gamma_*(\widetilde{\mathcal{M}}))^{\sim}
  \xrightarrow{\beta}\widetilde{\mathcal{M}}
\tag{3.3.5.1}
\]
\[
\label{II.3.3.5.2-fr}
  \Gamma_*(\mathcal{F})\xrightarrow{\alpha}
  \Gamma_*((\Gamma_*(\mathcal{F}))^{\sim})
  \xrightarrow{\Gamma_*(\beta)}\Gamma_*(\mathcal{F})
\tag{3.3.5.2}
\]
sont les isomorphismes identiques.
\end{proposition}

\begin{proof}
La question est en effet locale sur $Y$, et on est ramené à
(\hyperref[II.2.6.5-fr]{2.6.5}).
\end{proof}

\subsection{Conditions de finitude.}
\label{subsection:II.3.4-fr}

